\documentclass[11pt]{article}
\usepackage[margin=.7in]{geometry}\usepackage[T1]{fontenc}\usepackage{lmodern}\usepackage{xcolor}\usepackage{hyperref}
\definecolor{navy}{HTML}{102A43}\definecolor{blue}{HTML}{2F80ED}\definecolor{green}{HTML}{16866A}\definecolor{gold}{HTML}{9A6400}\hypersetup{colorlinks=true,linkcolor=blue,urlcolor=blue}\setlength{\parindent}{0pt}\setlength{\parskip}{4pt}
\newcommand{\slide}[1]{\clearpage{\Large\color{navy}\textbf{#1}}\par\vspace{4pt}}\newcommand{\field}[2]{\textbf{\color{blue}#1:} #2\par}
\title{\DeckTitle — Instructor Notes}\author{Computing Commons}\date{}
\begin{document}\sloppy\setlength{\emergencystretch}{3em}\maketitle
\section*{Study before class}\textbf{Day:} \DeckDay · \textbf{Week:} \DeckWeek\par\textbf{Exact source:} \DeckSource\par\textbf{Essential question:} \DeckQuestion\par\vspace{4pt}Read the exact canonical source named above and use the session README for its pinned path and commit. Review the most important concepts before class.
\slide{1. \DeckTitle — frame the move}\field{Purpose}{Make the essential question and the bounded finish line visible.}\field{Say}{\DeckPromise}\field{Ask}{What would a careful person need to know before accepting the first answer?}\field{Do}{Give students the one-sentence problem and let them predict.}\field{Watch for}{A polished answer being treated as evidence.}\field{Move on when}{Students can state the question and one uncertainty.}
\slide{2. Today's finish line}\field{Purpose}{Turn the lesson into an observable student move.}\field{Say}{\DeckMove}\field{Ask}{Which part is a claim, and which part is an observation?}\field{Do}{Point to the three columns and name the receipt.}\field{Watch for}{Students trying to complete every possible extension.}\field{Move on when}{Each student can name a bounded output.}
\slide{3. The mental model}\field{Purpose}{Give students a usable model without flattening the source.}\field{Say}{\DeckModel}\field{Ask}{Where does this model stop being enough?}\field{Do}{Use the boundary card explicitly.}\field{Watch for}{A metaphor replacing the actual criteria or evidence.}\field{Move on when}{Students can explain the model in their own words.}
\slide{4. Worked computing example}\field{Purpose}{Translate the idea into a safe, clearly labeled computing case.}\field{Say}{This application is instructor-created; it is not a claim from the primary source.}\field{Ask}{What would we inspect first?}\field{Do}{Have pairs mark the checkpoint and alternative explanation.}\field{Watch for}{Students copying the example instead of transferring the move.}\field{Move on when}{Pairs can defend one checkpoint.}
\slide{5. Ask before the answer}\field{Purpose}{Fade the scaffold by eliciting a prediction before revealing the frame.}\field{Say}{\DeckQuestion}\field{Ask}{What else could explain the same result?}\field{Do}{Collect one claim, one reason, and one uncertainty.}\field{Watch for}{Debate over conclusions without shared evidence.}\field{Move on when}{The room has at least two plausible explanations or criteria.}
\slide{6. Do this}\field{Purpose}{Rehearse the professional action, not just the vocabulary.}\field{Say}{\DeckDo}\field{Ask}{Which step makes the work easier to inspect?}\field{Do}{Time a short individual attempt, then a partner check.}\field{Watch for}{Scope creep or a hidden leap from observation to conclusion.}\field{Move on when}{A partner can reproduce the route from the receipt.}
\slide{7. Source and limits}\field{Purpose}{Protect source fidelity and bounded claims.}\field{Say}{\DeckSupports}\field{Ask}{What does this not prove?}\field{Do}{Read the limitation card aloud and invite one counterexample.}\field{Watch for}{Turning a pattern into a universal law or causal claim.}\field{Move on when}{Students can state one support and one limit.}
\slide{8. Recovery}\field{Purpose}{Make revision and help-seeking part of the method.}\field{Say}{A failed check is information about the route, not a verdict about identity.}\field{Ask}{What assumption should we inspect first?}\field{Do}{Model one bounded retry and preserve the original observation.}\field{Watch for}{Repeated retries that change several variables at once.}\field{Move on when}{Students name a smaller next attempt.}
\slide{9. Receipt}\field{Purpose}{Close with durable evidence and a next action.}\field{Say}{\DeckEvidence}\field{Ask}{How will another person know what you actually checked?}\field{Do}{Give students the four-part receipt frame.}\field{Watch for}{Narration without the observation or limit.}\field{Move on when}{Each student has a concise receipt.}
\slide{10. Exit}\field{Purpose}{Transfer the move to the next professional situation.}\field{Say}{\DeckExit}\field{Ask}{Where could this habit prevent a bad decision or unclear claim?}\field{Do}{Collect the exit sentence or point students to the next route.}\field{Watch for}{A new slogan without a concrete application.}\field{Move on when}{Students name a context and a verification step.}
\end{document}
